\documentclass[10pt,a4paper]{amsart}
\usepackage[margin=19mm]{geometry}
\usepackage{amsmath,amssymb,mathtools}
\usepackage[scr=rsfs,cal=euler]{mathalfa}
\usepackage{xfrac}
\usepackage[unicode,hidelinks,bookmarks=false]{hyperref}
\newcommand{\ie}{i.e.\ }
\newcommand{\cf}{cf.\ }
\newcommand{\resp}{respectively\ }
\newcommand{\Subsection}{\subsection}
\providecommand{\nobreakdash}{\nobreak\hbox{-}}
\setlength{\emergencystretch}{2em}
\multlinegap0pt
\usepackage{bm}
\usepackage{bbm}
\usepackage{dsfont}
\usepackage{xspace}
\usepackage{cancel}
\usepackage{mathtools}
\usepackage{amsthm}
\makeatletter
\@ifundefined{lemma}{\newtheorem{lemma}{Lemma}}{}
\makeatother
\DeclareMathOperator{\rank}{rank}
\title[A Slice-Rank Drift Bound for Random Quantum \(k\)-SAT]{A Slice-Rank Drift Bound for Random Quantum \(k\)-SAT}
\author{Jean Bernoulli Ravelomanana}
\date{Source: arXiv:2607.23847v1 (CC BY 4.0)}
\begin{document}
\maketitle

\noindent\emph{Source and licence.} A Slice-Rank Drift Bound for Random Quantum \(k\)-SAT. Version arXiv:2607.23847v1 (CC BY 4.0), distributed under the Creative Commons Attribution 4.0 International licence, as recorded in the arXiv licence metadata for this exact version and shown on the abstract page https://arxiv.org/abs/2607.23847v1. Full rights notice: licenses/arxiv-2607.23847-CC-BY-4.0.txt.\par\smallskip

\noindent\textit{Editorial scope.} Counted unit: the Lemma whose source label is \texttt{lem:multiplicative-slice-rank} in the article (source0.tex lines 693--707, source label \texttt{lem:multiplicative-slice-rank}), delivered together with the article's own complete original proof of it (source0.tex lines 709--938, one \texttt{proof} environment of 230 lines). No mathematics is written, repaired or re-derived: statement and proof are the article's own text line for line. Required context retained uncounted: none. Every definition, display and label of those lines is retained unchanged. Omitted from this package and not redistributed: 1--692, 708--708, 939--1807. Nothing inside a retained excerpt is omitted or altered: every retained body is delivered exactly as the article's lines are written, so no omission receipt of any kind accompanies this package. Citations: the retained excerpts carry no citation command at all, so no bibliography entry of the article is transcribed and no reference is redistributed. Reference adaptation: none was needed and none was made; every reference inside the delivered unit resolves inside it, so no unresolved reference or citation is produced. Printed slips kept literal: no printed word, symbol, hypothesis or number of the counted statement or of its proof was altered. Layout and class: the article's own class is \texttt{article} with packages including geometry, amsmath, amssymb, amsthm, mathtools, microtype, hyperref, cite; the wrapper is the lane's pinned \texttt{amsart} editorial wrapper at 10pt on A4, which restates the theorem-like environments the retained text needs and adds the article's own macro definitions that the retained text needs (\texttt{\textbackslash rank}), with extra wrapper packages (bm, bbm, dsfont, xspace, cancel, mathtools). The delivered numbering is the unit's own. Assets: no retained span contains a figure environment, a graphics-inclusion command, a PDF-page inclusion, a table or a TikZ picture; no image file is copied into this package, the package declares no asset dependency and no image file is redistributed with it. Licence: version arXiv:2607.23847v1 of this article is distributed under the Creative Commons Attribution 4.0 International licence, as recorded in the arXiv licence metadata for this exact version and shown on the abstract page https://arxiv.org/abs/2607.23847v1. A full rights notice is delivered at licenses/arxiv-2607.23847-CC-BY-4.0.txt. Characters: every retained excerpt is pure ASCII, so the delivered engine, XeLaTeX with its default Latin Modern text font, reports no missing character.
\begin{lemma}[Multiplicative slice-rank inequality]
	\label{lem:multiplicative-slice-rank}
	Let
	\[
	W\subseteq(\mathbb C^2)^{\otimes m}
	\]
	be a nonzero subspace, and set
	\[
	D=\dim W.
	\]
	Then
	\[
	\prod_{i=1}^m r_i(W)\geq D^{m-1}.
	\]
\end{lemma}

\begin{proof}
	We prove the result by induction on \(m\).
	
	For \(m=1\), the claim is
	\[
	r_1(W)\geq D^0=1,
	\]
	which holds because \(W\neq0\), and a generic covector on \(\mathbb C^2\)
	does not annihilate all of \(W\).
	
	Assume the claim known for \((m-1)\)-fold tensor products.  After relabelling
	the coordinates, split off the \(m\)-th tensor factor and write
	\[
	(\mathbb C^2)^{\otimes m}
	=
	\mathbb C^2\otimes\mathcal K,
	\qquad
	\mathcal K=(\mathbb C^2)^{\otimes(m-1)}.
	\]
	Let
	\[
	r=r_m(W).
	\]
	Choose a covector on the first factor whose contraction rank on \(W\) is \(r\).
	After a change of basis in \(\mathbb C^2\), we may assume that this covector
	is \(\langle0|\).  Let
	\[
	\pi_0:W\to\mathcal K
	\]
	be the corresponding \(|0\rangle\)-slice map.  Then
	\[
	\rank \pi_0=r.
	\]
	Set
	\[
	A=\pi_0(W)\subseteq\mathcal K,
	\qquad
	\dim A=r.
	\]
	The kernel of \(\pi_0|_W\) consists of vectors of the form
	\[
	|1\rangle\otimes b,
	\]
	where \(b\) belongs to some subspace
	\[
	B\subseteq\mathcal K.
	\]
	Hence
	\[
	\dim B=D-r.
	\]
	
	The two slice maps of \(W\) onto the \(|0\rangle\)- and
	\(|1\rangle\)-components both have rank at most \(r\), since \(r\) is the
	maximal contraction rank along the \(m\)-th coordinate.  Together, these two
	slices embed \(W\) into \(\mathcal K\oplus\mathcal K\).  Therefore
	\[
	D\leq 2r,
	\]
	and hence
	\[
	r\geq \frac D2.
	\]
	
	Choose a linear section of \(\pi_0:W\to A\).  Then there is a linear map
	\[
	F:A\to\mathcal K
	\]
	such that every element of \(W\) can be written uniquely in the form
	\[
	|0\rangle\otimes a
	+
	|1\rangle\otimes(Fa+b),
	\qquad
	a\in A,\ b\in B.
	\]
	
	Now fix a coordinate \(i<m\), viewed as a coordinate of \(\mathcal K\).  Let
	\[
	r_i(A)
	\qquad\text{and}\qquad
	r_i(B)
	\]
	be the corresponding generic contraction ranks of \(A\) and \(B\).  Choose a
	covector on the \(i\)-th coordinate of \(\mathcal K\) which is generic
	simultaneously for \(A\) and \(B\).  This is possible because the intersection
	of finitely many nonempty Zariski-open sets is again nonempty and Zariski-open.
	
	Let
	\[
	C_i:\mathcal K\to(\mathbb C^2)^{\otimes([m-1]\setminus\{i\})}
	\]
	be the associated contraction.  The contracted image of \(W\) along coordinate
	\(i\) is
	\[
	\left\{
	|0\rangle\otimes C_i a
	+
	|1\rangle\otimes(C_iFa+C_i b)
	:
	a\in A,\ b\in B
	\right\}.
	\]
	Projection onto the \(|0\rangle\)-component has image \(C_i(A)\), of dimension
	\(r_i(A)\).  Its kernel contains the subspace
	\[
	\left\{
	|1\rangle\otimes C_i b:b\in B
	\right\},
	\]
	which has dimension \(r_i(B)\).  Hence this contracted image has dimension at
	least
	\[
	r_i(A)+r_i(B).
	\]
	Since \(r_i(W)\) is the maximal contraction rank along coordinate \(i\), we
	obtain
	\[
	r_i(W)\geq r_i(A)+r_i(B)
	\qquad\text{for every } i<m.
	\]
	
	By the induction hypothesis applied to \(A\subseteq\mathcal K\),
	\[
	\prod_{i=1}^{m-1} r_i(A)\geq r^{m-2}.
	\]
	If \(B=0\), then \(D=r\), and therefore
	\[
	\prod_{i=1}^m r_i(W)
	\geq
	r\prod_{i=1}^{m-1}r_i(A)
	\geq
	D\cdot D^{m-2}
	=
	D^{m-1}.
	\]
	Thus we may assume \(B\neq0\).  Applying the induction hypothesis to
	\(B\subseteq\mathcal K\) gives
	\[
	\prod_{i=1}^{m-1} r_i(B)\geq (D-r)^{m-2}.
	\]
	
	Using \(r_i(W)\geq r_i(A)+r_i(B)\), we get
	\[
	\prod_{i=1}^{m-1} r_i(W)
	\geq
	\prod_{i=1}^{m-1}\bigl(r_i(A)+r_i(B)\bigr).
	\]
	We shall use the multiplicative Minkowski inequality
	\[
	\left(
	\prod_{i=1}^{m-1}(x_i+y_i)
	\right)^{1/(m-1)}
	\geq
	\left(
	\prod_{i=1}^{m-1}x_i
	\right)^{1/(m-1)}
	+
	\left(
	\prod_{i=1}^{m-1}y_i
	\right)^{1/(m-1)}
	\]
	for nonnegative \(x_i,y_i\).  This follows from the concavity and homogeneity
	of the geometric mean.  Therefore
	\[
	\prod_{i=1}^{m-1} r_i(W)
	\geq
	\left(
	r^{(m-2)/(m-1)}
	+
	(D-r)^{(m-2)/(m-1)}
	\right)^{m-1}.
	\]
	Multiplying by \(r_m(W)=r\), we obtain
	\[
	\prod_{i=1}^m r_i(W)
	\geq
	r
	\left(
	r^{(m-2)/(m-1)}
	+
	(D-r)^{(m-2)/(m-1)}
	\right)^{m-1}.
	\]
	
	It remains to check that the right-hand side is at least \(D^{m-1}\).  Set
	\[
	x=\frac rD.
	\]
	Since \(r\geq D/2\), we have \(x\in[1/2,1]\).  Dividing by \(D^{m-1}\), it
	suffices to prove
	\[
	x
	\left(
	x^{(m-2)/(m-1)}
	+
	(1-x)^{(m-2)/(m-1)}
	\right)^{m-1}
	\geq 1.
	\]
	Let
	\[
	a=x^{1/(m-1)},
	\qquad
	b=(1-x)^{1/(m-1)}.
	\]
	Then \(a\geq b\) and
	\[
	a^{m-1}+b^{m-1}=1.
	\]
	The left-hand side above, raised to the power \(1/(m-1)\), is
	\[
	a(a^{m-2}+b^{m-2})
	=
	a^{m-1}+ab^{m-2}.
	\]
	Since \(a\geq b\),
	\[
	ab^{m-2}\geq b^{m-1}.
	\]
	Thus
	\[
	a^{m-1}+ab^{m-2}
	\geq
	a^{m-1}+b^{m-1}
	=
	1.
	\]
	This proves the scalar inequality and completes the induction.
\end{proof}

\end{document}
